\begin{circuitikz}[font=\small]
% Společná regulovaná větev 5 V, samostatné přívody k zátěžím.
\node[anchor=west,font=\footnotesize] at (0,8.8) {Externí zdroj 5 V / alespoň 4 A};
\draw[thick] (0,7.5) node[left]{+5 V}
 to[fuse,l={F1 T 3{,}75 A}] (3,7.5)
 --(6,7.5)--(14,7.5);
\draw[thick,fill=light] (14,1.3) rectangle (16,8.1);
\node[align=center] at (15,5.4) {Servo\\MG996R};
\node[anchor=west] at (14.1,7.5) {V+};
\node[anchor=west] at (14.1,3.4) {SIG};
\node[anchor=west] at (14.1,1.9) {GND};
% Blokování u konektoru; značka GND zde výslovně označuje síť G0.
\draw (10,7.5) to[C] (10,5.5) node[ground]{};
\node[anchor=east,align=right] at (9.6,6.45) {C9\\1000 µF};
\node at (10.25,6.85) {$+$};
\node[below,font=\footnotesize] at (10,4.85) {G0};
\draw (12.5,7.5) to[C] (12.5,5.5) node[ground]{};
\node[anchor=east,align=right] at (12.1,6.45) {C10\\100 nF};
\node[below,font=\footnotesize] at (12.5,4.85) {G0};
% Arduino: vstupy vlevo, signál vpravo, zem dole.
\draw[thick,fill=light] (4.8,2) rectangle (7.3,5.8);
\node[align=center] at (6.05,4.6) {Arduino\\Uno R3};
\node[anchor=west] at (4.9,5.3) {5 V};
\node[anchor=west] at (4.9,3.4) {A0};
\node[anchor=east] at (7.2,3.4) {D9};
\node[above] at (6.05,2.1) {GND};
\draw (4.2,7.5)--(4.2,5.3);
\fill (4.2,7.5) circle(1.5pt) (4.2,5.3) circle(1.5pt);
\node[above,font=\footnotesize] at (6.05,5.9) {USB odpojeno};
% P2 jako dělič, jezdec W skutečně připojen přes R8 na A0.
\draw (4.8,5.3)--(1,5.3)--(1,3.9);
\draw (.8,2.9) rectangle (1.2,3.9);
\node[anchor=east,align=right] at (.6,3.4) {P2\\10 k$\Omega$};
\draw (1,2.9)--(1,1)--(6.05,1)--(6.05,2);
\draw[-{Latex[length=2mm]}] (1.9,3.4)--(1.2,3.4);
\node[above] at (1.6,3.4) {W};
\draw (1.9,3.4) to[R] (3.7,3.4)--(4.8,3.4);
\node[align=center] at (2.8,4.15) {R8\\1 k$\Omega$};
\draw (4,3.4) to[C,l2_={C11 and 100 nF}] (4,1);
% Signál je veden vodorovně bez křížení s blokovacími kondenzátory.
\draw (7.3,3.4) to[R] (9.4,3.4)--(14,3.4);
\node[align=center] at (8.35,4.15) {R6\\330 $\Omega$};
\draw (10,3.4) to[R,l2_={R7 and 100 k$\Omega$}] (10,0);
% Výkonový návrat a samostatná zem řídicí části se setkají v G0.
\draw[very thick] (14,1.9)--(13.4,1.9)--(13.4,0)--(0,0) node[left]{Zdroj $-$};
\draw (6.05,1)--(7.6,1)--(7.6,0);
\fill (7.6,0) circle(2.5pt);
\node[below] at (7.6,0) {G0};
\foreach \x/\y in {4/1,4/3.4,6.05/1,10/0,10/3.4,10/7.5,12.5/7.5}
 {\fill (\x,\y) circle(1.5pt);}
\end{circuitikz}
